\documentclass[11pt]{article}
\usepackage[margin=1in]{geometry}

% ------------------------------------------------------------
% Typography / encoding (prevents garbled characters like "-¿")
% ------------------------------------------------------------
\usepackage[T1]{fontenc}
\usepackage[utf8]{inputenc} % safe for pdflatex; omit if using lualatex/xelatex
\usepackage{lmodern}
\usepackage{microtype}
\setlength{\emergencystretch}{3em} % reduces overfull lines in narrow spots

% ------------------------------------------------------------
% Paths/URLs that can wrap (fixes long \texttt{...} overflow)
% ------------------------------------------------------------
\usepackage{xurl} % provides breakable \url/\path

% ------------------------------------------------------------
% Tables / lists
% ------------------------------------------------------------
\usepackage{booktabs}
\usepackage{array}
\newcolumntype{L}[1]{>{\raggedright\arraybackslash}p{#1}}
\usepackage{enumitem}
\setlist[itemize]{itemsep=0.2em, topsep=0.25em}

% ------------------------------------------------------------
% Prompts: highlighted monospace blocks with line-wrapping
% ------------------------------------------------------------
\usepackage{xcolor}
\usepackage{listings}
\usepackage{tcolorbox}
\tcbuselibrary{listings,skins,breakable}

\lstdefinestyle{promptstyle}{
  basicstyle=\ttfamily\small,
  columns=fullflexible,
  breaklines=true,
  breakatwhitespace=false,
  keepspaces=true,
  showstringspaces=false,
  postbreak=\mbox{\textcolor{gray}{$\hookrightarrow$}\space}
}

\newtcblisting{promptbox}{
  listing only,
  listing options={style=promptstyle},
  colback=gray!6,
  colframe=gray!45,
  boxrule=0.6pt,
  arc=2pt,
  left=6pt,right=6pt,top=6pt,bottom=6pt,
  enhanced,
  breakable
}

% ------------------------------------------------------------
% Convenience macros / task numbering
% ------------------------------------------------------------
\newcommand{\datasetfile}[1]{\path{#1}}

\newcounter{task}
\newcommand{\tasksection}[1]{%
  \refstepcounter{task}%
  \subsection*{Task \thetask: #1}%
}

\newenvironment{taskmeta}{%
  \begin{description}[
      style=nextline,
      leftmargin=3.0cm,
      labelwidth=2.7cm,
      itemsep=0.15em,
      topsep=0.25em
  ]%
}{\end{description}}

\begin{document}

\section*{Task Definitions and Prompts}

This section documents the benchmark's 16 tasks in a paper-ready format: dataset construction, model input/output, reward function, and the exact prompt template used.
Train/validation counts are taken from \datasetfile{dataset/minerva_base_split/metadata.json}.
Placeholders shown in braces (e.g., \texttt{\{CVE\_DESCRIPTION\}}) are substituted with instance-specific content at runtime.

% ============================================================
% 1--3: Mapping-explorer CVE -> ATT&CK
% ============================================================

\tasksection{CVE $\rightarrow$ ATT\&CK (Exploitation Technique)}
\textbf{Goal.} Map a vulnerability description to the ATT\&CK technique that best represents the \emph{exploitation method} used to trigger or leverage the vulnerability.
\begin{taskmeta}
  \item[Source] \datasetfile{dataset/mappings-explorer/kev-02.13.2025_attack-15.1-enterprise.yaml}. CVE descriptions are extracted from YAML comments (CVE IDs normalized to ``This vulnerability''/``this vulnerability''); labels are derived from \texttt{attack\_object\_id}.
  \item[Input] CVE description.
  \item[Output] One MITRE ATT\&CK Enterprise technique ID (T#### or T####.###).
  \item[Reward] \texttt{reward\_technique\_detection\_sigma}.
\end{taskmeta}
\begin{promptbox}
Given the vulnerability description below, output the single most appropriate MITRE ATT&CK Enterprise technique ID for:

Exploitation Technique - the method (technique) used to exploit the vulnerability.

Requirements:
- Use MITRE ATT&CK Enterprise technique IDs only.
- Return exactly ONE technique ID.
- Prefer a sub-technique ID (e.g., T1059.003) if it is more suitable; otherwise return the parent technique ID (e.g., T1059).

Vulnerability description:
{CVE_DESCRIPTION}
\end{promptbox}

\tasksection{CVE $\rightarrow$ ATT\&CK (Primary Impact)}
\textbf{Goal.} Predict the ATT\&CK technique that best captures the \emph{initial benefit} obtained immediately after exploitation.
\begin{taskmeta}
  \item[Source] Same mapping-explorer build as Task 1.
  \item[Input] CVE description.
  \item[Output] One ATT\&CK Enterprise technique ID.
  \item[Reward] \texttt{reward\_technique\_id}.
\end{taskmeta}
\begin{promptbox}
Given the vulnerability description below, output the single most appropriate MITRE ATT&CK Enterprise technique ID for:

Primary Impact - the initial benefit (impact) gained through exploitation of the vulnerability.

Requirements:
- Use MITRE ATT&CK Enterprise technique IDs only.
- Return exactly ONE technique ID.
- Prefer a sub-technique ID (e.g., T1059.003) if it is more suitable; otherwise return the parent technique ID (e.g., T1059).

Vulnerability description:
{CVE_DESCRIPTION}
\end{promptbox}

\tasksection{CVE $\rightarrow$ ATT\&CK (Secondary Impact)}
\textbf{Goal.} Given a vulnerability description and the primary impact already achieved, predict the next most likely adversary action enabled by that access/impact.
\begin{taskmeta}
  \item[Source] Same mapping-explorer build as Task 1.
  \item[Input] CVE description + primary impact technique ID.
  \item[Output] One ATT\&CK Enterprise technique ID.
  \item[Reward] \texttt{reward\_technique\_id}.
\end{taskmeta}
\begin{promptbox}
Given the vulnerability description and the primary impact already obtained, output the single most appropriate MITRE ATT&CK Enterprise technique ID for:

Secondary Impact - what the adversary can do by gaining the benefit of the primary impact.

Requirements:
- Use MITRE ATT&CK Enterprise technique IDs only.
- Return exactly ONE technique ID.
- Prefer a sub-technique ID (e.g., T1059.003) if it is more suitable; otherwise return the parent technique ID (e.g., T1059).

Vulnerability description:
{CVE_DESCRIPTION}

Primary impact:
{PRIMARY_IMPACT_TECHNIQUE_ID}
\end{promptbox}

% ============================================================
% 4--5: Sigma -> ATT&CK
% ============================================================

\tasksection{Sigma $\rightarrow$ ATT\&CK (Technique)}
\textbf{Goal.} Given a Sigma rule excerpt (log source + detection logic), predict the single ATT\&CK technique that the rule is intended to detect.
\begin{taskmeta}
  \item[Source] Sigma rules from \datasetfile{dataset/sigma/rules} and \datasetfile{dataset/sigma/rules-threat-hunting}. Each instance includes an excerpt containing (at minimum) title, \texttt{logsource}, and \texttt{detection}.
  \item[Input] Sigma rule excerpt.
  \item[Output] One ATT\&CK Enterprise technique ID.
  \item[Reward] \texttt{reward\_technique\_id}.
\end{taskmeta}
\begin{promptbox}
Given the Sigma rule excerpt below (log source + detection logic), provide the single most appropriate MITRE ATT&CK Enterprise technique ID that best represents the adversary behavior this rule is intended to detect.

Technique - how an adversary achieves a tactical objective by performing an action.

Requirements:
- Use MITRE ATT&CK Enterprise technique IDs only.
- Return exactly ONE technique ID.
- Prefer a sub-technique ID (e.g., T1059.003) if it is more suitable than the parent technique; otherwise return the parent technique ID (e.g., T1059).

Sigma rule excerpt:
{SIGMA_RULE_EXCERPT}
\end{promptbox}

\tasksection{Sigma $\rightarrow$ ATT\&CK (Tactics)}
\textbf{Goal.} Predict all ATT\&CK tactics (objectives) that are valid for the behavior targeted by the Sigma rule excerpt.
\begin{taskmeta}
  \item[Source] Same Sigma build as Task 4.
  \item[Input] Sigma rule excerpt.
  \item[Output] A set of ATT\&CK tactic IDs (TA000x).
  \item[Reward] \texttt{reward\_tactic\_ids}.
\end{taskmeta}
\begin{promptbox}
Given the Sigma rule excerpt below (log source + detection logic), enumerate all MITRE ATT&CK Enterprise tactic IDs (TA000x) that are valid for the adversary behavior this rule is intended to detect.

Tactic - why an adversary performs an action (the adversary's tactical objective).

Requirements:
- Use MITRE ATT&CK Enterprise tactic IDs (TA000x) only.
- List ALL applicable tactic IDs (multiple may apply).

Sigma rule excerpt:
{SIGMA_RULE_EXCERPT}
\end{promptbox}

% ============================================================
% 6--9: Scenario -> ATT&CK (MITRE procedures)
% ============================================================

\tasksection{Scenario $\rightarrow$ ATT\&CK (Technique)}
\textbf{Goal.} Map an ATT\&CK procedure-style scenario to the single best matching ATT\&CK technique.
\begin{taskmeta}
  \item[Source] Procedure scenarios derived from MITRE ATT\&CK procedure descriptions. Entity names are generalized by type (intrusion-set -> \texttt{A threat actor}, campaign -> \texttt{A campaign}, malware/tool -> \texttt{A malware}) while preserving leading articles.
  \item[Input] Adversary procedure scenario text.
  \item[Output] One ATT\&CK Enterprise technique ID.
  \item[Reward] \texttt{reward\_technique\_id}.
\end{taskmeta}
\begin{promptbox}
Given the adversary procedure description below, provide the single most appropriate MITRE ATT&CK Enterprise technique ID that best represents the behavior.

Requirements:
- Use MITRE ATT&CK Enterprise technique IDs only.
- Return exactly ONE technique ID.
- Prefer a sub-technique ID (e.g., T1059.003) if it is more suitable; otherwise return the parent technique ID (e.g., T1059).

Adversary procedure:
{SCENARIO_TEXT}
\end{promptbox}

\tasksection{Scenario $\rightarrow$ ATT\&CK (Tactics)}
\textbf{Goal.} Given a procedure scenario, select the \emph{most appropriate} ATT\&CK tactics that apply (with task-specific cardinality).
\begin{taskmeta}
  \item[Source] Same procedure scenario build as Task 6.
  \item[Input] Adversary procedure scenario text.
  \item[Output] Exactly \texttt{\{COUNT\}} tactic IDs (TA000x).
  \item[Reward] \texttt{reward\_tactic\_ids}.
\end{taskmeta}
\begin{promptbox}
Given the adversary procedure description below, list EXACTLY {COUNT} MITRE ATT&CK Enterprise tactic ID{S_SUFFIX} (TA000x) that apply to this behavior.

Requirements:
- Use MITRE ATT&CK Enterprise tactic IDs (TA000x) only.
- List EXACTLY {COUNT} tactic ID{S_SUFFIX} that are most appropriate.

Adversary procedure:
{SCENARIO_TEXT}
\end{promptbox}

\tasksection{Scenario $\rightarrow$ ATT\&CK (Mitigations)}
\textbf{Goal.} Predict the most relevant mitigations for the described adversary behavior (with task-specific cardinality).
\begin{taskmeta}
  \item[Source] Same procedure scenario build as Task 6.
  \item[Input] Adversary procedure scenario text.
  \item[Output] Exactly \texttt{\{COUNT\}} mitigation IDs (M####).
  \item[Reward] \texttt{reward\_mitigation\_ids}.
\end{taskmeta}
\begin{promptbox}
Given the adversary procedure description below, list EXACTLY {COUNT} MITRE ATT&CK Enterprise mitigation ID{S_SUFFIX} (M####) that best mitigate this behavior.

Requirements:
- Use MITRE ATT&CK Enterprise mitigation IDs only.
- List EXACTLY {COUNT} mitigation ID{S_SUFFIX} that are most appropriate.

Adversary procedure:
{SCENARIO_TEXT}
\end{promptbox}

% ============================================================
% 10--12: CVE -> CWE / CVSS (NVD)
% ============================================================

\tasksection{CVE $\rightarrow$ CWE (NVD)}
\textbf{Goal.} Infer the most likely weakness categories (CWE) from a CVE natural-language description.
\begin{taskmeta}
  \item[Source] Built from NVD CVE records.
  \item[Input] CVE description.
  \item[Output] A list of CWE IDs in \texttt{CWE-<number>} format (cardinality fixed by \texttt{\{COUNT\}}).
  \item[Reward] \texttt{reward\_cwe\_ids}.
\end{taskmeta}
\begin{promptbox}
Given the Common Vulnerabilities and Exposures (CVE) description below, list EXACTLY {COUNT} Common Weakness Enumeration (CWE) ID{S_SUFFIX} in CWE-<number> format.

CVE description:
{CVE_DESCRIPTION}
\end{promptbox}

\tasksection{CVE $\rightarrow$ CVSS v3.1 Base Vector (NVD)}
\textbf{Goal.} Predict the CVSS v3.1 base vector string directly from a CVE description.
\begin{taskmeta}
  \item[Source] Built from NVD CVE records with CVSS v3.1 scoring.
  \item[Input] CVE description.
  \item[Output] One CVSS v3.1 base vector string.
  \item[Reward] \texttt{reward\_cvss\_v31}.
\end{taskmeta}
\begin{promptbox}
Given the Common Vulnerabilities and Exposures (CVE) description below, provide the CVSS v3.1 base vector string (format: CVSS:3.1/AV:X/AC:X/PR:X/UI:X/S:X/C:X/I:X/A:X).

CVE description:
{CVE_DESCRIPTION}
\end{promptbox}

% ============================================================
% 13--15: CAPEC example tasks
% ============================================================

\tasksection{CAPEC Example $\rightarrow$ CAPEC ID}
\textbf{Goal.} Given an example attack description, classify it to the best matching CAPEC pattern ID.
\begin{taskmeta}
  \item[Source] Built from CAPEC example instances. (When multiple parent/child mappings exist, parent fallback is used for ATT\&CK mappings elsewhere.)
  \item[Input] Example attack description.
  \item[Output] A single CAPEC ID in \texttt{CAPEC-<number>} format.
  \item[Reward] \texttt{reward\_capec\_id}.
\end{taskmeta}
\begin{promptbox}
Given the example attack description below, provide the single best matching CAPEC ID in CAPEC-<number> format.

Example description:
{EXAMPLE_TEXT}
\end{promptbox}

\tasksection{CAPEC Example $\rightarrow$ CWE IDs}
\textbf{Goal.} Infer the most relevant CWE weakness categories from a CAPEC example description (limited to a small set per instance).
\begin{taskmeta}
  \item[Source] CAPEC example instances; CWE labels limited to at most 3 per example.
  \item[Input] Example attack description.
  \item[Output] Exactly \texttt{\{COUNT\}} CWE IDs in \texttt{CWE-<number>} format.
  \item[Reward] \texttt{reward\_cwe\_ids}.
\end{taskmeta}
\begin{promptbox}
Given the example attack description below, list EXACTLY {COUNT} Common Weakness Enumeration (CWE) ID{S_SUFFIX} in CWE-<number> format that apply.

Example description:
{EXAMPLE_TEXT}
\end{promptbox}

\tasksection{CAPEC Example $\rightarrow$ ATT\&CK Technique ID}
\textbf{Goal.} When an example maps cleanly to a single ATT\&CK technique, predict that technique from the example text.
\begin{taskmeta}
  \item[Source] CAPEC examples; ATT\&CK labels are included only when a single unambiguous mapping exists.
  \item[Input] Example attack description.
  \item[Output] One ATT\&CK Enterprise technique ID.
  \item[Reward] \texttt{reward\_technique\_id}.
\end{taskmeta}
\begin{promptbox}
Given the example attack description below, provide the single most appropriate MITRE ATT&CK Enterprise technique ID that best represents the adversary behavior.

Technique - how an adversary achieves a tactical objective by performing an action.

Requirements:
- Use MITRE ATT&CK Enterprise technique IDs only.
- Return exactly ONE technique ID.
- Prefer a sub-technique ID (e.g., T1059.003) if it is more suitable; otherwise return the parent technique ID (e.g., T1059).

Example description:
{EXAMPLE_TEXT}
\end{promptbox}

% ============================================================
% 16: Threat actor (procedures)
% ============================================================

\tasksection{Threat Actor (Procedures)}
\textbf{Goal.} Attribute an observed set of procedures to the most likely threat actor.
\begin{taskmeta}
  \item[Source] Built from MITRE ATT\&CK intrusion-set procedure descriptions. For each actor, we sample 60--100\% of techniques (min 3) and generate a number of questions based on technique-count bins chosen to balance total questions per bin ($\sim$100 each across 6 bins). Each bin assigns 3--8 questions per actor (bin 0 $\rightarrow$ 3, bin 5 $\rightarrow$ 8); the last bin is open-ended.
  \item[Input] A list of observed procedures.
  \item[Output] Threat actor name.
  \item[Reward] \texttt{reward\_threat\_actor\_name}.
  \item[Lookup] \texttt{threat\_actor\_lookup.json} is emitted alongside the dataset and copied into \texttt{rlvr/mydata/minerva\_base} (canonical names + aliases from MITRE ATT\&CK + supplemental aliases from \texttt{rlvr/verl/utils/reward\_score/aliases.csv}).
\end{taskmeta}
\begin{promptbox}
Given the observed adversary procedures below, identify the most likely threat actor.

Observed procedures:
{PROCEDURE_LIST}

Return only the threat actor name.
\end{promptbox}

% ============================================================
% Dataset statistics
% ============================================================

\section*{Dataset Statistics}

Table~\ref{tab:dataset-splits} summarizes the split sizes for each task.

\begin{table}[t]
\centering
\small
\caption{Dataset splits by task (train/validation).}
\label{tab:dataset-splits}
\begin{tabular}{@{}r L{0.26\textwidth} L{0.17\textwidth} L{0.21\textwidth} L{0.16\textwidth} r@{}}
\toprule
\# & File & Input & Output & Reward function & Train/Val \\
\midrule
1  & \datasetfile{cve_to_attack_exploitation}      & CVE description                  & ATT\&CK technique ID     & \texttt{reward\_technique\_id} & 245 / 20 \\
2  & \datasetfile{cve_to_attack_primary_impact}    & CVE description                  & ATT\&CK technique ID     & \texttt{reward\_technique\_id} & 210 / 20 \\
3  & \datasetfile{cve_to_attack_secondary_impact}  & CVE description                  & ATT\&CK technique ID     & \texttt{reward\_technique\_id} & 64 / 10 \\
4  & \datasetfile{sigma_to_attack_tactics}         & Sigma detection + logsource      & ATT\&CK tactic IDs       & \texttt{reward\_tactic\_ids}   & 950 / 50 \\
5  & \datasetfile{sigma_to_attack_technique}       & Sigma detection + logsource      & ATT\&CK technique ID     & \texttt{reward\_technique\_detection\_sigma} & 950 / 50 \\
6  & \datasetfile{art_to_attack_technique}         & ART procedure snippet            & ATT\&CK technique ID     & \texttt{reward\_technique\_detection\_art} & 950 / 50 \\
7  & \datasetfile{sentinel_to_attack_technique}    & Sentinel analytics rule          & ATT\&CK technique ID     & \texttt{reward\_technique\_detection\_sentinel} & 950 / 50 \\
8  & \datasetfile{splunk_to_attack_technique}      & Splunk detection rule            & ATT\&CK technique ID     & \texttt{reward\_technique\_detection\_splunk} & 280 / 20 \\
9  & \datasetfile{scenario_to_technique}           & ATT\&CK procedure scenario       & ATT\&CK technique ID     & \texttt{reward\_technique\_id} & 7780 / 220 \\
10 & \datasetfile{scenario_to_tactics}             & ATT\&CK procedure scenario       & ATT\&CK tactic IDs       & \texttt{reward\_tactic\_ids}   & 1950 / 50 \\
11 & \datasetfile{scenario_to_mitigations}         & ATT\&CK procedure scenario       & ATT\&CK mitigation IDs   & \texttt{reward\_mitigation\_ids} & 7780 / 220 \\
12 & \datasetfile{cve_to_cwe}                      & CVE description                  & CWE IDs                  & \texttt{reward\_cwe\_ids}      & 6696 / 220 \\
13 & \datasetfile{cve_to_cvss_v31}                 & CVE description                  & CVSS v3.1 vector         & \texttt{reward\_cvss\_v31}     & 1900 / 100 \\
14 & \datasetfile{threat_actor}                    & Observed procedures              & Threat actor name        & \texttt{reward\_threat\_actor\_name} & 779 / 60 \\
15 & \datasetfile{capec_example_to_capec}          & CAPEC example                    & CAPEC ID                 & \texttt{reward\_capec\_id}     & 340 / 40 \\
16 & \datasetfile{capec_example_to_cwe}            & CAPEC example                    & CWE IDs                  & \texttt{reward\_cwe\_ids}      & 176 / 20 \\
\midrule
\multicolumn{5}{@{}r}{\textbf{Total}} & \textbf{32000 / 1200} \\
\bottomrule
\end{tabular}
\end{table}

% ============================================================
% Reward notes
% ============================================================

\section*{Reward Function Notes}

\begin{description}[
  style=nextline,
  leftmargin=3.6cm,
  labelwidth=3.3cm,
  itemsep=0.2em,
  topsep=0.25em
]
  \item[\texttt{reward\_technique\_id}] Score is 1.0 if technique and sub-technique match exactly; 0.5 if the parent technique matches but the sub-technique differs or is omitted; otherwise 0.
  \item[\texttt{reward\_tactic\_ids}] F1 score over the predicted vs.\ ground-truth sets of tactic IDs.
  \item[\texttt{reward\_cwe\_ids}] F1 score over the predicted vs.\ ground-truth sets of CWE IDs.
\item[\texttt{reward\_cvss\_v31}] Compute CVSS v3.1 score distance (1 - \lvert truth - pred \rvert / 10); invalid vectors score 0. AthenaBench \texttt{athena-cti-vsp} uses delta=7.7.
  \item[\texttt{reward\_mitigation\_ids}] F1 score over the predicted vs.\ ground-truth sets of mitigation IDs.
  \item[Training-only short-reasoning penalty] If a boxed answer exists and the text before it has fewer than 20 word tokens (regex word count), subtract 0.05 and clamp at $\ge 0$.
  \item[\texttt{binary\_id}] Score is 1.0 if the predicted ID matches the ground truth; otherwise 0.
  \item[\texttt{reward\_threat\_actor\_name}] Score is 1.0 if the prediction matches the actor name or an alias; otherwise 0.
\end{description}

\end{document}
